\documentclass[a4paper]{arrowhead}

\usepackage[yyyymmdd]{datetime}
\usepackage{etoolbox}
\usepackage[utf8]{inputenc}
\usepackage[T1]{fontenc}
\usepackage{multirow}
\usepackage{longtable}
\usepackage{tikz}
\usetikzlibrary{positioning,shapes.multipart,arrows.meta}

\renewcommand{\dateseparator}{-}
\setlength{\emergencystretch}{3em}
\lstset{language=HTTP,breakatwhitespace=false,columns=fullflexible}

%% Special references
\newcommand{\mref}[2][]{{\textcolor{ArrowheadPurple}{\hyperref[sec:model:#2]{\ifstrempty{#1}{#2}{#1}}}}}
\newcommand{\oref}[1]{{\textcolor{ArrowheadBlue}{\hyperref[sec:operations:#1]{#1}}}}
\newcommand{\pdef}[1]{{\textcolor{ArrowheadGrey}{#1\label{sec:model:primitives:#1}\label{sec:model:primitives:#1s}\label{sec:model:primitives:#1es}}}}
\newcommand{\pref}[1]{{\textcolor{ArrowheadGrey}{\hyperref[sec:model:primitives:#1]{#1}}}}

% as template_sd.tex; the input and output arguments are passed already formatted (\mref or \pref)
\newrobustcmd\osubsection[3]{
  \addtocounter{subsection}{1}
  \addcontentsline{toc}{subsection}{\protect\numberline{\thesubsection}operation \textcolor{ArrowheadBlue}{#1}}
  \subsection*{
    \thesubsection\quad
    operation
    \textcolor{ArrowheadBlue}{#1}
    (\notblank{#2}{#2}{})
    \notblank{#3}{: #3}{}
  }
  \label{sec:operations:#1}
}
\newrobustcmd\msubsection[2]{
  \addtocounter{subsection}{1}
  \addcontentsline{toc}{subsection}{\protect\numberline{\thesubsection}#1 \textcolor{ArrowheadPurple}{#2}}
  \subsection*{\thesubsection\quad#1 \textcolor{ArrowheadPurple}{#2}}
  \label{sec:model:#2} \label{sec:model:#2s} \label{sec:model:#2es}
}
%%

\begin{document}
\sloppy

%% Arrowhead Document Properties
\ArrowheadTitle{certificateAuthorityManagement}
\ArrowheadServiceID{certificateAuthorityManagement}
\ArrowheadType{Service Description}
\ArrowheadTypeShort{SD}
\ArrowheadVersion{5.2.2}
\ArrowheadDate{2026-10-30}
\ArrowheadAuthor{Jerker Delsing}
\ArrowheadStatus{RELEASE CANDIDATE}
\ArrowheadContact{jerker.delsing@ltu.se}
\ArrowheadFooter{\href{www.arrowhead.eu}{www.arrowhead.eu}}
\ArrowheadSetup
%%

%% Front Page
\begin{center}
  \vspace*{1cm}
  \huge{\arrowtitle}

  \vspace*{0.2cm}
  \LARGE{\arrowtype}
  \vspace*{1cm}

  \Large{Service ID: \textit{"\arrowid"}}

  \vspace*{\fill}

  % Front Page Image

  \vspace*{1cm}
  \vspace*{\fill}

  % Front Page Abstract
  \begin{abstract}
    This document describes the \texttt{certificate\allowbreak{}Authority\allowbreak{}Management} service, version 1.0.0, produced by the CertificateAuthority support system of Arrowhead 5.2.2.
  \end{abstract}

  \vspace*{1cm}
\end{center}

\newpage
%%

%% Table of Contents
\tableofcontents
\newpage
%%

\section{Overview}
\label{sec:overview}

This document describes the \texttt{certificate\allowbreak{}Authority\allowbreak{}Management} service, which provides operator management of the certificate register and of trusted keys, including revocation.
It is produced by the CertificateAuthority system \cite{cau-sysd} in service version 1.0.0.

The rest of this document is organized as follows.
In Section \ref{sec:operations}, we describe the abstract message operations provided by the service.
In Section \ref{sec:model}, we end the document by presenting the data types used by the mentioned operations.
The concrete interface is specified in the HTTP/TLS/JSON IDD \cite{cau-certificateAuthorityManagement-http-idd}.

\subsection{Significant Prior Art}

Successor of the Arrowhead 4.6 Certificate Authority management endpoints.

\subsection{How This Service Is Meant to Be Used}

Sysop lists issued certificates, revokes them by record or by content (also certificates the CA did not issue through \texttt{sign}, e.g. deployment certificates), and maintains trusted keys. Onboarding may revoke the certificates it requested.

\subsection{Important Delimitations and Dependencies}
\label{sec:delimitations}

Only Sysop and Certificate\allowbreak{}Authority (Onboarding for its own records). Revocation cannot be undone; every change republishes the CRL.

\newpage

\section{Service Interface}
\label{sec:operations}

This section describes the interfaces to the \texttt{certificate\allowbreak{}Authority\allowbreak{}Management} service.
In particular, each subsection names an abstract operation, an input type and an output type, in that order.
The input type is named inside parentheses, while the output type is preceded by a colon.
Input and output types are only denoted when accepted or returned, respectively, by the interface in question.

All abstract data types named in this section are defined in Section \ref{sec:model}.

\begin{figure}[ht!]
  \centering
  \begin{tikzpicture}[font=\small]
    \node[draw=ArrowheadBlue, very thick, rectangle split, rectangle split parts=2, align=center,
          rectangle split part fill={ArrowheadBlue!12,white}, text width=9cm, inner sep=6pt] (b)
      {\guillemotleft{}Service\guillemotright{}\\\textbf{certificateAuthorityManagement}
        \nodepart{second} \texttt{list-\allowbreak{}certificates}\\ \texttt{revoke-\allowbreak{}certificate}\\ \texttt{revoke-\allowbreak{}certificate-\allowbreak{}by-\allowbreak{}content}\\ \texttt{list-\allowbreak{}keys}\\ \texttt{add-\allowbreak{}key}\\ \texttt{delete-\allowbreak{}key}};
  \end{tikzpicture}
  \caption{SysML-style block of the \texttt{certificate\allowbreak{}Authority\allowbreak{}Management} service and its operations.}
  \label{fig:SD}
\end{figure}

The following interface operations are available.

\osubsection{list-certificates}{}{\pref{List}$<$\allowbreak{}\mref{CertificateRecord}$>$}

Returns the certificate register, newest first.

\osubsection{revoke-certificate}{\pref{Identifier}}{\mref{RemovalResult}}

Revokes a register record by identifier (\texttt{\{id\}}) and republishes the CRL.

\osubsection{revoke-certificate-by-content}{\mref{RevocationRequest}}{\mref{RevocationResult}}

Revokes any certificate of this CA given by content; unknown ones are added to the register as revoked.

\osubsection{list-keys}{}{\pref{List}$<$\allowbreak{}\mref{TrustedKeyRecord}$>$}

Returns the trusted keys.

\osubsection{add-key}{\mref{TrustedKeyRequest}}{\mref{TrustedKeyResult}}

Adds a trusted key with validity.

\osubsection{delete-key}{\pref{Identifier}}{\mref{RemovalResult}}

Deletes a trusted key by identifier (\texttt{\{id\}}).

\newpage

\section{Information Model}
\label{sec:model}

Here, all data objects that can be part of the \texttt{certificate\allowbreak{}Authority\allowbreak{}Management} service are listed in alphabetic order.
Note that each subsection, which describes one type of object, begins with the \textit{struct} keyword, which is used to denote a collection of named fields, each with its own data type.
As a complement to the explicitly defined types in this section, there is also a list of implicit primitive types in Section \ref{sec:model:primitives}, which are used to represent things like hashes and identifiers.

\msubsection{struct}{CertificateRecord}




\noindent\begin{tabularx}{\textwidth}{| >{\raggedright\arraybackslash}p{3.9cm} | >{\raggedright\arraybackslash\small}p{5.3cm} | X |} \hline
\rowcolor{gray!33} Field & Type & Description \\ \hline
id & \pref{Number} & Record identifier. \\ \hline
serial\_number & \pref{String} & Serial number. \\ \hline
common\_name & \pref{String} & Subject common name. \\ \hline
created\_by & \pref{String} & Requesting system (\texttt{revoked-\allowbreak{}by:\allowbreak{}\textless{}caller\textgreater{}} for imported revocations). \\ \hline
valid\_from & \pref{String} & Start of validity. \\ \hline
valid\_until & \pref{String} & End of validity. \\ \hline
revoked\_at & \pref{String} & Revocation time, null if not revoked. \\ \hline
certificate\_pem & \pref{String} & The certificate, PEM. \\ \hline
\end{tabularx}



\msubsection{struct}{RemovalResult}




\noindent\begin{tabularx}{\textwidth}{| >{\raggedright\arraybackslash}p{3.9cm} | >{\raggedright\arraybackslash\small}p{5.3cm} | X |} \hline
\rowcolor{gray!33} Field & Type & Description \\ \hline
success & \pref{Boolean} & Whether a record was changed; 404 when false. \\ \hline
\end{tabularx}



\msubsection{struct}{RevocationRequest}




\noindent\begin{tabularx}{\textwidth}{| >{\raggedright\arraybackslash}p{3.9cm} | >{\raggedright\arraybackslash\small}p{5.3cm} | X |} \hline
\rowcolor{gray!33} Field & Type & Description \\ \hline
certificate & \pref{String} & Certificate to revoke, PEM or Base64 DER; must be issued by this cloud CA. \\ \hline
\end{tabularx}



\msubsection{struct}{RevocationResult}




\noindent\begin{tabularx}{\textwidth}{| >{\raggedright\arraybackslash}p{3.9cm} | >{\raggedright\arraybackslash\small}p{5.3cm} | X |} \hline
\rowcolor{gray!33} Field & Type & Description \\ \hline
id & \pref{Number} & Register record of the certificate. \\ \hline
serial\allowbreak{}Number & \pref{String} & Serial number. \\ \hline
revoked & \pref{Boolean} & Always true. \\ \hline
changed & \pref{Boolean} & False if it was already revoked. \\ \hline
\end{tabularx}



\msubsection{struct}{TrustedKeyRecord}




\noindent\begin{tabularx}{\textwidth}{| >{\raggedright\arraybackslash}p{3.9cm} | >{\raggedright\arraybackslash\small}p{5.3cm} | X |} \hline
\rowcolor{gray!33} Field & Type & Description \\ \hline
id & \pref{Number} & Record identifier. \\ \hline
public\_key\_pem & \pref{String} & The key, PEM. \\ \hline
description & \pref{String} & Description. \\ \hline
valid\_after & \pref{String} & Start of validity. \\ \hline
valid\_before & \pref{String} & End of validity. \\ \hline
created\_at & \pref{String} & Creation time. \\ \hline
\end{tabularx}



\msubsection{struct}{TrustedKeyRequest}




\noindent\begin{tabularx}{\textwidth}{| >{\raggedright\arraybackslash}p{3.9cm} | >{\raggedright\arraybackslash\small}p{5.3cm} | X |} \hline
\rowcolor{gray!33} Field & Type & Description \\ \hline
public\allowbreak{}Key & \pref{String} & PEM Subject\allowbreak{}Public\allowbreak{}Key\allowbreak{}Info (alias \texttt{publicKeyPem}). \\ \hline
description & \pref{String} & Description. \\ \hline
valid\allowbreak{}After & \pref{DateTime} & Start of validity; default now. \\ \hline
valid\allowbreak{}Before & \pref{DateTime} & End of validity; default now plus \texttt{CERT\_\allowbreak{}VALIDITY\_\allowbreak{}DAYS}. \\ \hline
\end{tabularx}



\msubsection{struct}{TrustedKeyResult}




\noindent\begin{tabularx}{\textwidth}{| >{\raggedright\arraybackslash}p{3.9cm} | >{\raggedright\arraybackslash\small}p{5.3cm} | X |} \hline
\rowcolor{gray!33} Field & Type & Description \\ \hline
id & \pref{Number} & Record identifier. \\ \hline
trusted & \pref{Boolean} & Always true. \\ \hline
\end{tabularx}



\subsection{Primitives}
\label{sec:model:primitives}

Types and structures mentioned throughout this document that are assumed to be available to implementations of this service.
The concrete interpretations of each of these types and structures must be provided by any IDD document claiming to implement this service.

\noindent\begin{tabularx}{\textwidth}{| p{4.6cm} | X |} \hline
\rowcolor{gray!33} Type & Description \\ \hline
\pdef{Boolean} & One out of \texttt{true} or \texttt{false}. \\ \hline
\pdef{DateTime} & Pinpoints a specific moment in time (UTC). \\ \hline
\pdef{Identifier} & A string that identifies a record, e.g. a UUID or a composite identifier. \\ \hline
\pdef{List} & An \textit{array} of a known number of items, each having type A. \\ \hline
\pdef{Number} & An integer number. \\ \hline
\pdef{String} & A chain of characters. \\ \hline
\end{tabularx}

\newpage

\bibliographystyle{IEEEtran}
\bibliography{bibliography}

\newpage

\section{Revision History}
\subsection{Amendments}

\noindent\begin{tabularx}{\textwidth}{| p{1cm} | p{3cm} | p{2cm} | X | p{4cm} |} \hline
\rowcolor{gray!33} No. & Date & Version & Subject of Amendments & Author \\ \hline

1 & 2026-10-22 & \arrowversion & Document created from the Arrowhead LaTeX templates for the 5.2.2 Java Spring implementation, including the local extensions of this workspace. & Jerker Delsing \\ \hline
2 & 2026-10-30 & \arrowversion & Checked against the primary Java source tree (Eclipse-Arrowhead-java): services, operations, paths and data models unchanged. & Claude Opus 5.5 (draft) \\ \hline
\end{tabularx}

\subsection{Quality Assurance}

\noindent\begin{tabularx}{\textwidth}{| p{1cm} | p{3cm} | p{2cm} | X |} \hline
\rowcolor{gray!33} No. & Date & Version & Approved by \\ \hline

1 & & \arrowversion  &  \\ \hline

\end{tabularx}

\end{document}
